%%%%%%%%%%%%%%%%%
% CV targeted for Space Data Operations & Processing Engineer
%%%%%%%%%%%%%%%%%

\documentclass[8pt,a4paper,ragged2e,withhyper]{../_shared/altacv}

\geometry{left=1.25cm,right=1.25cm,top=.5cm,bottom=1.cm,columnsep=1.2cm}

\usepackage{paracol}
\usepackage{fontawesome5}
\usepackage{hyperref}

\iftutex 
  \setmainfont{Roboto Slab}
  \setsansfont{Lato}
  \renewcommand{\familydefault}{\sfdefault}
\else
  \usepackage[rm]{roboto}
  \usepackage[defaultsans]{lato}
  \renewcommand{\familydefault}{\sfdefault}
\fi

\definecolor{SlateGrey}{HTML}{2E2E2E}
\definecolor{LightGrey}{HTML}{666666}
\definecolor{DarkPastelRed}{HTML}{8AC0D9}
\definecolor{PastelRed}{HTML}{00B0DA}
\definecolor{GoldenEarth}{HTML}{B4AD0C}

\colorlet{name}{black}
\colorlet{tagline}{PastelRed}
\colorlet{heading}{DarkPastelRed}
\colorlet{headingrule}{GoldenEarth}
\colorlet{subheading}{PastelRed}
\colorlet{accent}{PastelRed}
\colorlet{emphasis}{SlateGrey}
\colorlet{body}{LightGrey}

\definecolor{violet}{HTML}{5A2582}
\definecolor{rouge}{HTML}{F53B3B}
\definecolor{noir}{HTML}{00232C}
\definecolor{bleu}{HTML}{00B0DA}
\definecolor{rose}{HTML}{F831A0}
\definecolor{jaune}{HTML}{FFEC00}
\definecolor{vert}{HTML}{34AD6C}

\renewcommand{\namefont}{\Huge\rmfamily\bfseries}
\renewcommand{\personalinfofont}{\footnotesize}
\renewcommand{\cvsectionfont}{\LARGE\rmfamily\bfseries}
\renewcommand{\cvsubsectionfont}{\large\bfseries}
\renewcommand{\cvItemMarker}{{\small\textbullet}}
\renewcommand{\cvRatingMarker}{\faCircle}

\input{../_shared/pubs-num.tex}
\addbibresource{../_shared/sample.bib}

\begin{document}

\name{BOILEAU Guillaume}
\tagline{Space Data Operations \& Processing Engineer | Python, Linux, pipelines, monitoring \& anomaly analysis}

\photoL{2.8cm}{../_shared/moi}

\personalinfo{%
  \email{guillaume.boileaupro@gmail.com}
  \phone{+33 6 59 94 85 84}
  \location{Alpes-Maritimes, FRANCE}
  \linkedin{guillaume-boileau-phd-891b81265}
  \researchgate{Guillaume-Boileau-2}
  \orcid{0000-0002-3576-6968}
}

\makecvheader
\vspace{-0.3cm}

\AtBeginEnvironment{itemize}{\small}
\columnratio{0.64}

\begin{paracol}{2}

\cvsection{Experience}

\cvevent{\textbf{Software Engineer - System Integration, Data Flows \& Operational Monitoring}}{VEV Platform Services France}{2025 -- 2026}{Valbonne, France}

\textbf{Context:} Development and integration of software services for an EV charging CPMS platform connected to operational infrastructure, charging stations, data flows and hardware systems.

\textbf{Objective:} Support reliable operational data-flow continuity, monitoring, incident analysis and software/system integration in an industrial environment.

\begin{itemize}
  \item \textbf{Operational systems:} Contributed to the development, deployment support and maintenance of software services connected to operational infrastructure and field equipment.
  \item \textbf{Data flows:} Implemented, monitored and corrected FTP-based operational data exchange services involving structured data and recurring processing actions.
  \item \textbf{Monitoring \& reliability:} Used operational logs, OpenSearch and Sentry to investigate incidents, follow system behaviour and support service reliability.
  \item \textbf{Anomaly analysis:} Analysed issues involving software behaviour, operational data exchanges, interface continuity and hardware-connected systems.
  \item \textbf{Structured data:} Manipulated SQL, relational concepts and JSON-based objects/configurations for daily technical actions and workflows.
  \item \textbf{System integration:} Worked on interactions between platform services, connected charging stations, data exchanges, electrical circuits, communication buses and operational behaviours.
  \item \textbf{Technical reporting:} Documented incidents, technical findings, corrective actions and operational observations for internal teams.
  \item \textbf{Industrial environment:} Worked with software, operations and technical stakeholders to clarify issues, follow priorities and support iterative delivery.
  \item \textbf{Tooling:} Used Git-based workflows, Zenhub, Confluence, OpenSearch and Sentry for delivery follow-up, diagnostics and technical documentation.
\end{itemize}

\divider

\cvevent{\textbf{Senior R\&D Engineer - Space Data Processing, Validation \& Technical Coordination}}{CNES, Laboratoire Lagrange, Observatoire de la Côte d'Azur}{2023 -- 2025}{Nice, France}

\textbf{Context:} Engineering and data-analysis activities for the LISA ESA/NASA space mission, involving complex software chains, simulation outputs, data products, validation workflows and international coordination.

\textbf{Objective:} Develop, integrate and validate Python-based data processing workflows under strong consistency, traceability, performance and verification constraints.

\begin{itemize}
  \item \textbf{Space data processing:} Developed Python workflows for the processing, comparison and validation of complex simulation outputs and space-mission data products.
  \item \textbf{Pipeline development:} Developed CATpyLISA, a Python platform for large-scale model integration, comparison, validation and technical review.
  \textcolor{DarkPastelRed}{\href{https://gitlab.oca.eu/gboileau/lisa_l3_tools}{\bf GitLab: CATpyLISA}}
  \item \textbf{Operational logic:} Structured repeatable processing sequences, reruns, consistency checks, data-product consolidation and validation-oriented diagnostics.
  \item \textbf{Reprocessing workflows:} Supported iterative processing and re-analysis of heterogeneous outputs to assess robustness, consistency and impact of configuration changes.
  \item \textbf{Data validation:} Ensured traceability of inputs, configurations, outputs, assumptions, metadata and validation results.
  \item \textbf{Anomaly investigation:} Analysed discrepancies between models, outputs and expected behaviours to identify causes, limitations and corrective directions.
  \item \textbf{Performance analysis:} Analysed sensitivity to modelling assumptions, numerical parameters, data-processing choices and pipeline configurations.
  \item \textbf{Signal processing:} Analysed weak, noisy and superposed signals using spectral methods, signal/noise separation logic and performance metrics.
  \item \textbf{Technical reporting:} Produced analysis notes, validation summaries, technical documentation and review material for contributors and project stakeholders.
  \item \textbf{Coordination:} Coordinated contributors, supervised interns and supported technical activities across multiple stakeholders in an international environment.
\end{itemize}

\divider

\cvevent{\textbf{Data Scientist - Noise Diagnostics, Data Analysis \& Performance Validation}}{Universiteit Antwerpen, Belgium}{2021 -- 2023}{Antwerp, Belgium}

\textbf{Context:} Data analysis and performance studies for precision instruments in a high-requirement international environment linked to gravitational-wave detector performance.

\textbf{Objective:} Identify, characterize and mitigate factors affecting measurement quality and system sensitivity using statistical modelling, signal processing and validation-oriented methods.

\begin{itemize}
  \item \textbf{Data processing:} Analysed measurement and simulation outputs to identify abnormal behaviours, noise contributions and degradation sources.
  \item \textbf{Signal processing:} Analysed instrumental and environmental effects affecting measurement quality, sensitivity and stability.
  \item \textbf{Noise diagnostics:} Developed methods for spectral analysis, noise characterization, sensitivity studies and cross-sensor correlations.
  \item \textbf{Performance monitoring:} Identified contributors to performance degradation and supported prioritisation of correction or mitigation actions.
  \item \textbf{Reproducibility:} Produced traceable, reproducible and documented analyses for international scientific and engineering stakeholders.
  \item \textbf{Technical communication:} Communicated complex results through structured reporting, presentations and collaborative review.
\end{itemize}

\divider

\cvevent{\textbf{Junior R\&D Engineer - Simulation, Signal Analysis \& Scientific Computing}}{ARTEMIS, Observatoire de la Côte d'Azur}{2018 -- 2021}{Nice, France}

\textbf{Context:} Data-analysis and simulation activities for gravitational-wave mission preparation, involving weak-signal detection, modelling, simulation and noise characterization.

\textbf{Objective:} Develop robust analysis methods for complex signals and support technical investigations in demanding system environments.

\begin{itemize}
  \item \textbf{Simulation workflows:} Built simulation approaches for complex signal environments, uncertainty studies and large-scale datasets.
  \item \textbf{Signal analysis:} Developed methods for the separation of overlapping low-SNR signals in complex datasets.
  \item \textbf{Algorithm assessment:} Compared simulated outputs, model behaviours and analysis results to assess consistency, robustness and performance limitations.
  \item \textbf{Technical investigations:} Investigated noise sources through cross-sensor correlation analyses in diagnostic contexts.
  \item \textbf{Scientific computing:} Developed strong expertise in Python-based analysis, modelling, documentation and reproducible workflows.
\end{itemize}

\switchcolumn

\cvsection{Profile}

\textbf{Engineer at the interface between space data processing, operational workflows, software integration, validation and anomaly analysis}. Background developed through ESA/NASA space mission activities, scientific data pipelines, precision instrumentation and industrial operational systems.

Experienced in \textbf{Python}, \textbf{Linux}, \textbf{Shell/Bash}, \textbf{SQL}, \textbf{Docker}, structured data handling, JSON workflows, data processing pipelines, validation logic, monitoring, technical reporting and system reliability analysis.

For a \textbf{Space Data Operations \& Processing Engineer} role, I bring a strong engineering profile: developing and operating data workflows, analysing anomalies, validating data products, supporting reprocessing activities, documenting results and working across software, data, operations and system constraints.

\cvsection{Fit for Space Data Operations}

\begin{itemize}
  \item Strong experience in Python-based data processing pipelines for space-mission and scientific data products.
  \item Solid Linux, Bash, Git, Docker and reproducible workflow background.
  \item Experience with SQL, relational database concepts and daily manipulation of JSON-based structured data.
  \item Strong ability to analyse anomalies, investigate discrepancies and contribute to corrective actions.
  \item Experience with operational monitoring, logs, OpenSearch, Sentry and production-like incident analysis.
  \item Background in validation, traceability, technical documentation and performance-oriented reporting.
  \item Strong fit for international technical environments requiring autonomy, rigor, analytical thinking and technical English.
  \item Space-domain experience from LISA ESA/NASA; directly transferable to satellite data processing and operational data workflows.
\end{itemize}

\cvsection{Team Management}

\begin{itemize}
  \item Supervision of \textbf{3 Master's thesis interns (M2)} during software, modelling and data-analysis activities.
  \item Coordination of technical activities involving \textbf{research engineers and technicians} on a \textbf{data processing and validation pipeline} for the \textbf{LISA space mission}.
  \item Experience organizing tasks, supporting contributors, reviewing outputs and maintaining alignment across collaborative technical developments.
\end{itemize}

\cvsection{Languages}

\cvskill{English}{4}
\divider
\cvskill{Spanish}{2}
\divider

\medskip

\cvsection{Education}

\cvevent{PhD in Physics}{Université Côte d'Azur}{2018 -- 2021}{Nice, France}

\divider

\cvevent{Selective Graduate Program in Fundamental Physics (Magistère)}{Université Paris-Saclay}{2016 -- 2018}{Orsay, France}

\divider

\cvevent{MSc in Astronomy and Astrophysics}{Observatoire de Paris / Université Paris-Saclay}{2017 -- 2018}{Paris, France}

\divider

\cvevent{BSc in Fundamental Physics}{Université Paris-Saclay}{2015 -- 2016}{Orsay, France}

\end{paracol}

\cvsection{Professional Training}

\cvevent{Supervised Machine Learning (Regression \& Classification)}{DeepLearning.AI - Andrew Ng}{2020}{Coursera}

\divider

\cvevent{Recherche reproductible : principes méthodologiques pour une science transparente}{FUN MOOC -- Inria}{2025}{France Université Numérique}

\divider

\cvevent{Continuous Professional Training -- Software, Data, AI and Systems}{OpenClassrooms}{2024 -- 2026}{}

\small{
Python, SQL, Machine Learning, AI, Linux, JavaScript, TypeScript, Angular, Node.js, Django, REST APIs, C/C++, Java, algorithms, relational databases, software engineering fundamentals and data workflows.
}

\cvsection{Projects}

\cvevent{CNES / LISA -- Space Data Processing, Validation \& Reprocessing Workflows}{ESA/NASA Space Mission - Large-scale International Collaboration}{2023 -- 2025}{Nice, France}

Contribution to Python-based data processing, model-comparison and validation workflows for the LISA space mission within a large international engineering and scientific collaboration.

\textbf{Space data relevance:} Processing of heterogeneous simulation outputs, consolidation of data products, traceability of configurations, validation of results and technical review.

\textbf{Operations relevance:} Repeatable processing sequences, reruns, consistency checks, anomaly investigation, structured outputs and reporting for stakeholders.

\textbf{Critical systems relevance:} Mission-level performance analysis, validation evidence, rigorous documentation and work in a high-constraint international aerospace context.

\divider

\cvevent{Einstein Telescope / LIGO--Virgo--KAGRA -- Data Diagnostics \& Detector Performance}{International Collaborations}{2021 -- 2023}{Antwerp, Belgium}

Contribution to collaborative developments in data analysis, signal characterization, statistical diagnostics and performance investigations within large-scale international scientific and engineering collaborations.

\textbf{Role:} Development of analysis tools, contribution to technical studies, reporting and collaborative review in high-standard environments.

\textbf{System impact:} Studies on environmental and instrumental noise sources supporting the identification of performance limitations and design constraints for next-generation gravitational-wave observatories.

\divider

\cvevent{VEV -- Operational Data Flows, Monitoring \& System Integration}{EV Charging Infrastructure / CPMS Platform}{2025 -- 2026}{Valbonne, France}

Contribution to software development, operational diagnostics and hardware/software interface analysis for a CPMS platform connected to EV charging infrastructure.

\textbf{Operational relevance:} Monitoring of data exchanges, incident analysis, logs investigation, OpenSearch/Sentry usage, software reliability, JSON/SQL handling and operational traceability.

\cvsection{Strengths}

\vspace{-.3cm}

\cvsubsection{Space data processing \& operations}

{\small
\cvtag{Space Data Processing}
\cvtag{Satellite Data}
\cvtag{Data Operations}
\cvtag{Data Processing Pipelines}
\cvtag{Reprocessing Workflows}
\cvtag{Operational Workflows}
\newline
\cvtag{Pipeline Automation}
\cvtag{Data Validation}
\cvtag{Data Product Consolidation}
\cvtag{Processing Chains}
\cvtag{Workflow Reliability}
\cvtag{Operational Reporting}
\cvtag{Technical Documentation}
\cvtag{Technical English}
}

\divider\smallskip
\vspace{-.3cm}
\newline

\cvsubsection{Programming, Linux \& data tools}

{\small
\cvtag{Python}
\cvtag{Shell/Bash}
\cvtag{Linux}
\cvtag{SQL}
\cvtag{Relational Databases}
\cvtag{PostgreSQL Basics}
\cvtag{JSON}
\cvtag{Structured Data}
\cvtag{Pandas}
\cvtag{NumPy}
\cvtag{Scikit-Learn}
\cvtag{Jupyter}
\cvtag{Git}
\cvtag{GitLab}
\cvtag{GitHub}
\cvtag{Docker}
\cvtag{CI/CD}
\cvtag{REST API}
\cvtag{TypeScript}
\cvtag{Angular}
\cvtag{Node.js}
\cvtag{Java}
\cvtag{C}
\cvtag{C++}
}

\divider\smallskip
\vspace{-.3cm}

\cvsubsection{Monitoring, reliability \& anomaly analysis}

{\small
\cvtag{Operational Monitoring}
\cvtag{OpenSearch}
\cvtag{Sentry}
\cvtag{Log Analysis}
\cvtag{Incident Analysis}
\cvtag{Anomaly Analysis}
\cvtag{Corrective Actions}
\cvtag{Root-Cause Analysis}
\cvtag{System Reliability}
\cvtag{Performance Monitoring}
\cvtag{Data-Flow Monitoring}
\cvtag{Issue Investigation}
\cvtag{Robustness}
\cvtag{Problem Solving}
\cvtag{Analytical Thinking}
}

\divider\smallskip
\vspace{-.3cm}

\cvsubsection{Signal, diagnostics \& complex data}

{\small
\cvtag{Signal Processing}
\cvtag{Weak Signal Detection}
\cvtag{Low-SNR Analysis}
\cvtag{Noise Modelling}
\cvtag{Noise Diagnostics}
\cvtag{Spectral Analysis}
\cvtag{Cross-Sensor Correlation}
\cvtag{Simulation Workflows}
\cvtag{Synthetic Data}
\cvtag{Uncertainty Analysis}
\cvtag{Sensitivity Analysis}
\cvtag{Performance Assessment}
\cvtag{Complex Data Analysis}
\cvtag{Scientific Computing}
}

\divider\smallskip
\vspace{-.3cm}

\cvsubsection{Validation, quality \& critical systems}

{\small
\cvtag{System Validation}
\cvtag{IVVQ}
\cvtag{V-Cycle}
\cvtag{Verification}
\cvtag{Validation Evidence}
\cvtag{Traceability}
\cvtag{Requirements Awareness}
\cvtag{Technical Evidence}
\cvtag{Data Reliability}
\cvtag{Quality Checks}
\cvtag{Consistency Checks}
\cvtag{Risk Identification}
\cvtag{Critical Systems}
\cvtag{Aerospace Context}
\cvtag{Space Systems}
}

\divider\smallskip
\vspace{-.3cm}

\cvsubsection{Distributed computing, delivery \& collaboration}

{\small
\cvtag{HPC}
\cvtag{Slurm}
\cvtag{HTCondor}
\cvtag{Distributed Computing}
\cvtag{Containerized Environments}
\cvtag{Cloud Awareness}
\newline
\cvtag{Technical Workshops}
\cvtag{Technical Reporting}
\cvtag{Knowledge Sharing}
\cvtag{Stakeholder Alignment}
\cvtag{Cross-Functional Coordination}
\cvtag{Agile Environment}
\cvtag{Confluence}
\cvtag{JIRA}
\cvtag{Zenhub}
\cvtag{Autonomy}
\cvtag{Rigor}
\cvtag{Team Player}
}

\end{document}